\documentclass[mode=notes, palette=ocean, lang=french, fontsize=11pt]{modernclassnotes}

\course{Documentation du Paquet LaTeX}
\coursecode{MCN v1.2.0}
\professor{Équipe de Développement Modern Class Notes}
\institution{Open Source TeX Initiative}
\term{Édition 2026}
\docsubtitle{Manuel d'utilisation \& Guide d'architecture du paquet Class Notes}
\title{La classe \& paquet \texttt{modernclassnotes}}

\begin{document}

\maketitle

\tableofcontents
\newpage

\section{Introduction}

Le paquet et la classe \texttt{modernclassnotes} offrent un cadre LaTeX moderne, esthétique et open-source conçu spécialement pour les enseignants, professeurs et universitaires. Il permet de créer facilement des notes de cours, des fiches de synthèse et des feuilles d'exercices.

\begin{takeaways}
Points forts de \texttt{modernclassnotes} :
\begin{itemize}
    \item \textbf{Quatre modes spécialisés} : \texttt{notes} (polycopiés complets), \texttt{handout} (fiches résumées de 1--2 pages), \texttt{worksheet} (feuilles d'exercices avec corrigés) et \texttt{chapternotes} (résumés par chapitre).
    \item \textbf{Palettes de couleurs choisies} : Alternez entre 8 thèmes (\texttt{ocean}, \texttt{nord}, \texttt{midnight}, \texttt{emerald}, \texttt{burgundy}, \texttt{amethyst}, \texttt{purple}, \texttt{mono}) en une seule option.
    \item \textbf{Support multilingue natif} : Prise en charge avec \texttt{lang=english}, \texttt{lang=french}, \texttt{lang=german} et \texttt{lang=chinese}.
    \item \textbf{Zéro dépendance lourde} : Moteur de boîtes TeX/LaTeX natif assurant une portabilité maximale sur TeX Live, MiKTeX et Overleaf.
    \item \textbf{Suivi des cours \& chapitres} : Bannières de cours et de chapitres automatisées avec objectifs optionnels via \texttt{\textbackslash lecture}, \texttt{\textbackslash chapter} et \texttt{\textbackslash subchapter}.
\end{itemize}
\end{takeaways}

\section{Options de la classe \& Configuration}

Chargez la classe dans l'en-tête de votre document :
\begin{codebox}[Chargement de la classe]
\begin{lstlisting}[language=TeX]
\documentclass[mode=notes, palette=ocean, lang=french, fontsize=11pt]{modernclassnotes}
\end{lstlisting}
\end{codebox}

\subsection{Tableau des options}
\begin{table}[h]
\centering
\begin{tabular}{l|p{9.5cm}|l}
\toprule
\textbf{Clé} & \textbf{Valeurs} & \textbf{Défaut} \\
\midrule
\texttt{mode} & \texttt{notes}, \texttt{handout}, \texttt{worksheet}, \texttt{chapternotes} & \texttt{notes} \\ \hline
\texttt{palette} & \texttt{ocean}, \texttt{midnight}, \texttt{nord}, \texttt{emerald}, \texttt{burgundy}, \texttt{amethyst}, \texttt{purple}, \texttt{mono} & \texttt{ocean} \\ \hline
\texttt{fontsize} & \texttt{10pt}, \texttt{11pt}, \texttt{12pt} & \texttt{11pt} \\ \hline
\texttt{font} & \texttt{default}, \texttt{palatino}, \texttt{charter}, \texttt{sans}/\texttt{helvet}, \texttt{kpfonts}, \texttt{pagella}, \texttt{sourceserif}, \texttt{fira}, \texttt{libertine}, \texttt{utopia} & \texttt{default} \\ \hline
\texttt{lang} & \texttt{english} (\texttt{en}), \texttt{german} (\texttt{de}), \texttt{french} (\texttt{fr}), \texttt{chinese} (\texttt{zh}) & \texttt{english} \\ \hline
\texttt{official} & \texttt{true}, \texttt{false} & \texttt{false} \\ \hline
\texttt{showsolutions} & \texttt{true}, \texttt{false} & \texttt{true} \\ \hline
\texttt{parskip} & \texttt{true}, \texttt{false} & \texttt{true} \\
\bottomrule
\end{tabular}
\caption{Options de classe et paramètres}
\end{table}

\section{Commandes de métadonnées}

Définissez les métadonnées de votre cours dans le préambule :
\begin{codebox}[Configuration des métadonnées]
\begin{lstlisting}[language=TeX]
\course{Mécanique Quantique}
\coursecode{PHYS 301}
\professor{Prof. Claire Dubois}
\institution{Sorbonne Université}
\term{Semestre d'Automne 2026}
\docsubtitle{Fiche de synthèse – Cours 3}
\title{Équation de Schrödinger}
\end{lstlisting}
\end{codebox}

\section{Systèmes de suivi des cours et chapitres}

Utilisez la commande \texttt{\textbackslash lecture} pour diviser le document en séances :
\begin{codebox}[Macro de cours]
\begin{lstlisting}[language=TeX]
\lecture[12 Novembre 2026]{1}{Équation de Schrödinger et Puits de Potentiel}
\end{lstlisting}
\end{codebox}

Pour les notes structurées par chapitre (\texttt{mode=chapternotes}), utilisez \texttt{\textbackslash chapter} et \texttt{\textbackslash subchapter} :
\begin{codebox}[Macro de chapitre avec objectifs]
\begin{lstlisting}[language=TeX]
\chapter{1}{Systèmes Numériques et Nombres Binaires}{%
  Comprendre la numération binaire.,%
  Convertir entre les bases binaires et décimales.,%
  Former les compléments.%
}

\subchapter{2}{Nombres Binaires}
\end{lstlisting}
\end{codebox}

\section{Environnements pédagogiques}

\begin{definition}[Environnements d'encadrés pédagogiques]
\texttt{modernclassnotes} fournit des environnements structurés pour le contenu académique :
\begin{itemize}
    \item \texttt{theorem} (Théorème), \texttt{lemma} (Lemme), \texttt{proposition}, \texttt{corollary} (Corollaire), \texttt{proof} (Démonstration)
    \item \texttt{definition}, \texttt{example} (Exemple), \texttt{remark} (Remarque), \texttt{warning} (Attention), \texttt{tip} (Astuce), \texttt{takeaways}
    \item \texttt{exercise} (Exercice), \texttt{solution} (Solution), \texttt{codebox} (Extrait de code)
\end{itemize}
\end{definition}

\begin{example}[Exemple d'environnement Théorème]
\begin{lstlisting}[language=TeX]
\begin{theorem}[Identité d'Euler]
Pour tout nombre réel $\theta$, $e^{i\theta} = \cos\theta + i\sin\theta$.
\end{theorem}
\end{lstlisting}
\end{example}

\section{Directives open-source}

Pour redistribuer ce paquet :
\begin{enumerate}
    \item Conservez les termes de licence LPPL 1.3c / MIT dans le fichier \texttt{LICENSE}.
    \item Fournissez les fichiers PDF compilés dans le dossier \texttt{examples/}.
\end{enumerate}

\end{document}
